% ============================================================
%  第 3 章  系统需求分析与总体设计
% ============================================================

\chapter{系统需求分析与总体设计}

% ----------------------------------------------------------------
\section{功能需求分析}

根据开题报告任务书的要求及实际舆情监测场景调研，本系统需满足以下六类核心功能需求。

\subsection{多平台数据采集需求}

系统需能够对接多个国内主流社交媒体平台，以关键词为维度对相关话题下的公开评论数据进行采集。具体要求包括：支持微博、B 站、抖音、小红书、知乎、快手、贴吧等不少于五个主流平台；采集内容包含评论正文、发布时间、点赞数量、转发数量等关键字段；支持以轮询方式持续抓取最新内容，轮询间隔可配置；采集结果以结构化 JSON 格式输出，兼容 Kafka 消息协议。

\subsection{实时流数据接入需求}

系统需构建稳定的流数据消费管道，满足以下要求：以 Apache Kafka 为消息中间件，支持多生产者并发写入同一主题；Spark Structured Streaming 消费端能够在指定触发周期内完成对新增消息的处理；消费进度需持久化存储，支持异常重启后的断点续传；消息解析过程需对格式异常、字段缺失等情况进行容错处理。

\subsection{情感分析与热度计算需求}

系统需对每条评论进行情感极性判定与热度量化：情感分类需覆盖正向、负向、中性三个类别；情感推理过程需支持批量处理，以满足微批次场景下的吞吐要求；热度评分需综合考量点赞数、转发数与情感极性三个维度；风险等级需根据热度与情感标签的组合进行三级划定，为告警机制提供输入依据。

\subsection{数据存储与管理需求}

系统需提供可靠、可查询的数据持久化能力：原始数据、清洗数据与聚合指标分层存储，层间数据质量递进；存储层需支持 ACID 事务，防止并发写入导致的数据不一致；聚合指标需覆盖小时粒度与日粒度两个时间维度；支持话题维度的历史数据查询，数据范围至少覆盖近 365 天；每个话题下需支持按热度排名的 Top10 评论记录存储。

\subsection{负面舆情告警需求}

系统需具备自动化舆情预警能力：当某话题在某统计小时内负面评论数量超过最低阈值（默认 10 条）且负面占比超过预设阈值（默认 40\%）时，自动触发告警；告警通知以电子邮件形式发送，邮件内容包含话题名称、统计时段、评论数量、负面占比及高热度负面评论样例；同一话题同一统计小时内不应重复发送告警；告警事件需持久化存储，供后续查询与审计。

\subsection{可视化展示需求}

系统需提供直观、交互友好的数据展示界面：支持多话题切换，默认显示最新数据话题；展示话题核心 KPI 指标；提供评论数、负面评论数、负面占比、热度四类趋势图表，支持按天和按小时两种粒度切换；展示评论明细列表与历史告警记录；图表支持缩放与拖拽交互，适配不同时间跨度的数据展示需求。

% ----------------------------------------------------------------
\section{非功能需求分析}

\subsection{实时性需求}

系统端到端处理延迟（从评论写入 Kafka 到前端可查询结果更新）需满足准实时要求。在当前 60 秒微批次触发周期内，新增评论数据应在下一个完整批次处理完成后即可在前端呈现，整体端到端延迟控制在 120 秒以内。

\subsection{可靠性需求}

系统核心链路需具备容错恢复能力：Kafka 消费进度通过 Checkpoint 机制持久化，保证 Spark 任务异常重启后数据不丢失；Delta Lake 的 ACID 事务保证写入失败不会产生脏数据；告警去重机制保证不因重复处理产生告警轰炸。

\subsection{可扩展性需求}

系统架构需支持水平扩展：新增监测话题无需修改代码，仅需在爬虫端增加关键词配置即可自动纳入监测范围；新增社交媒体平台时，仅需在爬虫工厂中注册新的爬虫类，不影响下游处理逻辑。

\subsection{可维护性需求}

各模块职责清晰，接口定义规范：爬虫模块、消息队列、流处理管道、API 服务与前端之间均通过标准化接口（Kafka JSON Schema、REST API）进行解耦；日志输出规范，关键处理节点有明确的进度与错误日志，便于问题定位。

% ----------------------------------------------------------------
\section{系统总体架构设计}

\subsection{架构分层说明}

基于上述需求分析，本系统采用五层分布式架构设计，由下至上依次为：数据采集层、消息队列层、流计算层、数据存储层与展示层。系统整体架构如图~\ref{fig:architecture}~所示。

% ---- TikZ：五层架构图 ----
\begin{figure}[htbp]
\centering
\resizebox{\textwidth}{!}{%
\begin{tikzpicture}[
  node distance=0.55cm,
  arr/.style={-Stealth, thick, color=gray!65, shorten >=2pt, shorten <=2pt},
  lbl/.style={font=\small, color=gray!60, fill=white, inner sep=1pt},
]

% ---------- 各层节点（从上到下，用 below=of 定位）----------
\node[archlay, fill=violet!18]                    (L5)
  {\textbf{展示层（前端）}\\[3pt]
   \small Vue~3 + Vite \enspace·\enspace Apache ECharts 6\\
   多话题切换 \enspace·\enspace KPI 卡 \enspace·\enspace 趋势折线图 \enspace·\enspace 评论明细表 \enspace·\enspace 告警面板};

\node[archlay, fill=green!18, below=of L5]        (L4)
  {\textbf{数据服务层（API）}\\[3pt]
   \small FastAPI \enspace·\enspace 监听端口 18000\\
   \code{/api/dashboard} · \code{/api/comments} · \code{/api/alerts} 等 9 个 REST 端点};

\node[archlay, fill=orange!18, below=of L4]       (L3)
  {\textbf{数据存储层（湖仓）}\\[3pt]
   \small Delta Lake（本地文件系统）\\
   Bronze \enspace·\enspace Silver \enspace·\enspace Gold Hourly / Daily / Top10 \enspace·\enspace Alert Events};

\node[archlay, fill=blue!14, below=of L3]         (L2)
  {\textbf{流计算层（Spark Structured Streaming）}\\[3pt]
   \small 60\,s 微批次触发 \enspace·\enspace StructBERT 情感推理 \enspace·\enspace 热度计算\\
   \code{foreachBatch}: Bronze 写入 → Silver Merge → Gold 重算 → 告警检测};

\node[archlay, fill=cyan!14, below=of L2]         (L1)
  {\textbf{消息队列层（Apache Kafka）}\\[3pt]
   \small Topic: \code{public\_opinion\_raw} \quad Broker: \code{172.16.156.151:9092}\\
   持久化日志 \enspace·\enspace At-Least-Once 消费语义 \enspace·\enspace Offset 检查点管理};

\node[archlay, fill=teal!18, below=of L1]         (L0)
  {\textbf{数据采集层（爬虫）}\\[3pt]
   \small MediaCrawler 多平台异步爬虫 \enspace·\enspace CDP 反检测模式 \enspace·\enspace 轮询采集\\
   微博 \enspace·\enspace B 站 \enspace·\enspace 抖音 \enspace·\enspace 小红书 \enspace·\enspace 知乎 \enspace·\enspace 快手 \enspace·\enspace 贴吧};

% ---------- 层间箭头与标注 ----------
\draw[arr] (L0.north) -- (L1.south)
  node[lbl, midway, right=4mm] {JSONL / Kafka Producer};
\draw[arr] (L1.north) -- (L2.south)
  node[lbl, midway, right=4mm] {readStream（消息流）};
\draw[arr] (L2.north) -- (L3.south)
  node[lbl, midway, right=4mm] {foreachBatch 写入};
\draw[arr] (L3.north) -- (L4.south)
  node[lbl, midway, right=4mm] {deltalake 直接读取};
\draw[arr] (L4.north) -- (L5.south)
  node[lbl, midway, right=4mm] {HTTP REST / JSON};

\end{tikzpicture}}%
\caption{系统五层分布式架构示意图}
\label{fig:architecture}
\end{figure}

各层职责简述如下：\textbf{数据采集层}以爬虫程序对接多平台，将评论数据序列化为 JSON 写入本地 JSONL 文件，由 Kafka Producer 进程读取并推送至消息主题；\textbf{消息队列层}以 Kafka 承担生产者与消费者之间的解耦与缓冲，保障高吞吐、持久化与可重放；\textbf{流计算层}以 Spark Structured Streaming 60~秒微批次为周期消费 Kafka 消息，在 Driver 端完成情感推理与热度计算，并通过 \code{foreachBatch} 将结果写入 Delta Lake；\textbf{数据存储层}采用 Delta Lake 湖仓架构，按 Bronze-Silver-Gold 三层渐进式精炼数据；\textbf{数据服务层}由 FastAPI 提供 RESTful 接口，直接读取 Delta Lake 数据并返回 JSON；\textbf{展示层}由 Vue~3 + ECharts 构成前端可视化大屏。

\subsection{技术架构关键决策}

\paragraph{决策一：微批次模式 vs. 纯流式处理}
本系统选择 Spark Structured Streaming 的微批次（Micro-Batch）模式，而非 Apache Flink 的原生流式（Native Streaming）处理。主要考量如下：其一，StructBERT 模型推理在 Driver 端以批量方式执行，微批次模式天然匹配批量推理的计算节奏；其二，Delta Lake 的 Merge 操作在批次粒度下效率显著优于逐条写入；其三，PySpark 与 Python NLP 生态（ModelScope）的结合在成熟度与生产可用性上优于 PyFlink。

\paragraph{决策二：\texttt{foreachBatch} vs. 标准 Sink}
标准 Structured Streaming Sink（如 \code{delta} 格式的 \code{writeStream}）仅支持单目标写入，而本系统需同时向 Bronze、Silver、多个 Gold 表及 Alert 表写入。\code{foreachBatch} 接口允许在每个批次内执行完整的批处理逻辑，完美支持多目标扇出（Fan-Out）写入场景。

\paragraph{决策三：Delta Lake vs. 传统关系型数据库}
相较于将聚合结果写入 MySQL 或 PostgreSQL，Delta Lake 方案在存储成本、历史版本查询（Time Travel）以及流批一体三个维度上具有显著优势，更适合实时分析场景。

% ----------------------------------------------------------------
\section{数据流设计}

\subsection{数据采集阶段数据流}

爬虫进程将采集到的评论经字段归一化处理后，序列化为如表~\ref{tab:kafka-schema}~所定义的 JSON 格式写入 Kafka 主题。

\begin{table}[htbp]
\centering
\caption{Kafka 消息字段定义（\code{public\_opinion\_raw} 主题）}
\label{tab:kafka-schema}
\begin{tabular}{llll}
\toprule
\tablehead{字段名} & \tablehead{类型} & \tablehead{是否必填} & \tablehead{说明} \\
\midrule
\code{event\_id}     & string  & 必填 & 唯一事件标识，用于全系统去重 \\
\code{platform}      & string  & 可选 & 来源平台（bili / wb / xhs 等） \\
\code{topic}         & string  & 必填 & 话题关键词 \\
\code{publish\_time} & string  & 必填 & 评论发布时间（yyyy-MM-dd HH:mm:ss） \\
\code{time\_id}      & string  & 可选 & 时间桶标识 \\
\code{text}          & string  & 必填 & 评论正文 \\
\code{like\_count}   & long    & 可选 & 点赞数（缺省为 0） \\
\code{repost\_count} & long    & 可选 & 转发数（缺省为 0） \\
\code{ingest\_time}  & string  & 可选 & 数据入库时间戳 \\
\code{source\_file}  & string  & 可选 & 来源 JSONL 文件路径（溯源用） \\
\bottomrule
\end{tabular}
\end{table}

\subsection{流处理阶段数据流}

Spark Structured Streaming 从 Kafka 读取消息后，经过图~\ref{fig:pipeline}~所示的五步处理链路生成各层数据。

% ---- TikZ：批次处理流程图 ----
\begin{figure}[htbp]
\centering
\begin{tikzpicture}[
  node distance=0.5cm and 0.6cm,
  box/.style={
    draw=gray!60, rounded corners=5pt, thick,
    minimum width=2.4cm, minimum height=0.85cm,
    text width=2.2cm, align=center, font=\small
  },
  kafka/.style={box, fill=cyan!15},
  proc/.style={box, fill=blue!12},
  store/.style={box, fill=orange!15},
  alert/.style={box, fill=red!12},
  arr/.style={-Stealth, thick, gray!65},
  lbl/.style={font=\footnotesize, color=gray!60},
  stepnum/.style={font=\footnotesize, color=gray!55},
]

% ── 第一行：主处理链路 ──
\node[kafka]  (src)    {Kafka\\ Source};
\node[proc,   right=of src]    (parse)  {JSON\\ 解析};
\node[proc,   right=of parse]  (filter) {字段\\ 过滤};
\node[proc,   right=of filter] (bert)   {StructBERT\\ 推理};
\node[proc,   right=of bert]   (heat)   {热度\\ 计算};

% ── 步骤编号（各节点正上方，above=9pt 保证不压住边框）──
\node[stepnum, above=9pt of src]    {\textcircled{1}};
\node[stepnum, above=9pt of parse]  {\textcircled{2}};
\node[stepnum, above=9pt of filter] {\textcircled{3}};
\node[stepnum, above=9pt of bert]   {\textcircled{4}};
\node[stepnum, above=9pt of heat]   {\textcircled{5}};

% ── 第二行：写入链路（主链下方 1.6cm，与主链有充足间距）──
\node[store, below=1.6cm of parse]  (bronze) {Bronze\\ 写入};
\node[store, right=of bronze]       (silver) {Silver\\ Merge};
\node[store, right=of silver]       (gold)   {Gold\\ 重算};
\node[alert, right=of gold]         (alrt)   {告警\\ 检测};

% ── 主链箭头 ──
\draw[arr] (src)    -- (parse);
\draw[arr] (parse)  -- (filter);
\draw[arr] (filter) -- (bert);
\draw[arr] (bert)   -- (heat);

% ── 热度计算向下分支：先垂直下降到折角坐标，再折向 Bronze ──
\coordinate (corner) at ($(heat.south)+(0,-0.65)$);
\draw[arr] (heat.south) -- (corner) -| (bronze.north);
% enriched_df 标注在折角右侧，远离节点和箭头
\node[lbl, right=3pt] at (corner) {enriched\_df};

% ── 写入链路箭头 ──
\draw[arr] (bronze) -- (silver) node[lbl, midway, above=15pt] {追加};
\draw[arr] (silver) -- (gold)   node[lbl, midway, above=15pt] {Merge};
\draw[arr] (gold)   -- (alrt)   node[lbl, midway, above=15pt] {聚合};

% ── 告警旁路 ──
\draw[arr, dashed, red!50] (alrt.south) -- ++(0,-0.5)
  node[lbl, below=2pt, color=red!60] {SMTP};

\end{tikzpicture}
\caption{\code{foreachBatch} 批次处理五步流程图}
\label{fig:pipeline}
\end{figure}

各步骤说明如下：\textcircled{1}~Kafka Source 按偏移量读取最新消息（每批最多 50 条）；\textcircled{2}~通过 \code{from\_json} 按预定义 Schema 解析消息体，追加 Kafka 元数据字段；\textcircled{3}~过滤 \code{event\_id}、\code{text}、\code{publish\_time}、\code{topic} 任一为空的记录；\textcircled{4}~批量调用 StructBERT Pipeline，为每条记录追加五个情感字段；\textcircled{5}~计算 \code{heat\_score}、\code{risk\_level}、\code{event\_hour}、\code{event\_date} 等衍生字段，随后按 Bronze → Silver → Gold → Alert 顺序写入 Delta Lake。

\subsection{数据服务阶段数据流}

FastAPI 服务收到前端 HTTP 请求后，通过 \code{deltalake} Python 库（Delta Lake 纯 Python 实现，无需启动 JVM）直接读取 Delta Lake 表文件，转换为 Pandas DataFrame 进行过滤、排序与序列化处理，最终以 JSON 格式返回。整个读取过程平均响应时间在 1.3~秒以内（详见第~\ref{sec:api-perf}~节实测数据）。

% ----------------------------------------------------------------
\section{模块划分}
\label{sec:modules}

根据系统架构设计，将系统划分为八个功能子模块，如表~\ref{tab:modules}~所示。

\begin{table}[htbp]
\centering
\caption{系统功能子模块划分（M1--M8）}
\label{tab:modules}
\begin{tabular}{clp{4.5cm}p{5.1cm}}
\toprule
\tablehead{编号} & \tablehead{模块名称} & \tablehead{对应文件/目录} & \tablehead{主要职责} \\
\midrule
M1 & 多平台爬虫模块      & \code{crawler/}                          & 多平台评论采集、JSONL 输出 \\
M2 & Kafka 消费解析模块  & \tcode{stream\_single\_topic\_}\newline\tcode{to\_delta.py} & Kafka 接入、Schema 解析、字段过滤 \\
M3 & 情感分析服务模块    & \tcode{structbert\_sentiment.py}           & StructBERT 批量推理、标签归一化、异常降级 \\
M4 & 热度计算模块        & \tcode{build\_enriched\_df()}              & 热度公式计算、风险等级划定 \\
M5 & Delta Lake 写入模块 & \tcode{process\_batch()} /\newline\tcode{rebuild\_gold\_tables()} & Bronze/Silver/Gold 三层写入 \\
M6 & 告警检测通知模块    & \tcode{alert\_mailer.py} /\newline\tcode{check\_and\_send\_alerts()} & 阈值检测、邮件发送、告警持久化 \\
M7 & REST API 服务模块   & \code{api/app.py}                         & 9 个端点的数据查询与序列化 \\
M8 & 前端可视化模块      & \code{frontend/}                          & 大屏渲染、交互逻辑、图表更新 \\
\bottomrule
\end{tabular}
\end{table}

% ----------------------------------------------------------------
\section{关键接口定义}

\subsection{Kafka 消息接口}

Kafka 消息接口规范如下：Topic 名称为 \code{public\_opinion\_raw}；消息格式为 UTF-8 编码的 JSON 字符串；\code{event\_id}、\code{topic}、\code{publish\_time}、\code{text} 为必填字段；\code{like\_count} 与 \code{repost\_count} 缺省值为 0。

\subsection{REST API 接口规范}

所有 API 端点均为 HTTP GET 方法，响应格式为 JSON。核心端点规范如表~\ref{tab:api}~所示。

\begin{table}[htbp]
\centering
\caption{REST API 核心端点规范}
\label{tab:api}
\begin{tabular}{lp{4.5cm}p{5.5cm}}
\toprule
\tablehead{端点} & \tablehead{主要参数} & \tablehead{响应字段} \\
\midrule
\code{GET /api/health}     & 无 & 各 Delta 表版本与存在状态 \\
\code{GET /api/topics}     & 无 & 全部话题列表（按评论总量排序） \\
\code{GET /api/dashboard}  & \code{topic}、\code{daily\_limit}、\code{hourly\_limit} & \code{meta}、\code{summary}、\code{daily}、\code{hourly}、\code{alerts} \\
\code{GET /api/comments}   & \code{topic}、\code{view}、\code{date}、\code{hour}、\code{limit} & \code{rows}、\code{bucket\_options}、\code{data} \\
\code{GET /api/alerts}     & \code{topic}、\code{limit} & \code{rows}、\code{data} \\
\bottomrule
\end{tabular}
\end{table}

\subsection{告警触发规则接口}

告警触发条件通过环境变量配置，实现参数与代码的分离，如表~\ref{tab:alert-config}~所示。

\begin{table}[htbp]
\centering
\caption{告警相关环境变量配置}
\label{tab:alert-config}
\begin{tabular}{lll}
\toprule
\tablehead{环境变量名} & \tablehead{默认值} & \tablehead{含义} \\
\midrule
\code{ALERT\_MIN\_COMMENTS}          & \code{10}   & 触发告警的最低评论数阈值 \\
\code{ALERT\_NEG\_RATIO\_THRESHOLD}  & \code{0.40} & 触发告警的负面评论占比阈值 \\
\code{ALERT\_TOPN\_TEXTS}            & \code{3}    & 告警邮件中附带的高热度负面样例数 \\
\code{STRUCTBERT\_NEUTRAL\_THRESHOLD}& \code{0.72} & 中性标签判定的置信度阈值 \\
\code{STRUCTBERT\_BATCH\_SIZE}       & \code{32}   & 情感推理的批量大小 \\
\bottomrule
\end{tabular}
\end{table}
